% !TEX root = ../Tesis.tex
\chapter{Introducción general}
\section{Contexto general}
La contaminación del aire se define como la presencia de uno o más
contaminantes en la atmósfera en cantidades y duraciones que pueden
provocar daños al ecosistema. La composición normal del aire incluye
nitrógeno (78.09\%), oxígeno (20.95\%), argón (0.93\%), dióxido de
carbono (0.04\%), pequeñas cantidades de otros gases como helio, y
vapor de agua. A su vez, esta composición cuenta con isótopos que
pueden tener comportamientos diferentes en el sistema terrestre
(Qu \& Dan, 2024)

Los isótopos son átomos de un mismo elemento con diferente masa
molecular. El análisis isotópico de gases y aerosoles en la atmósfera
muestra diferencias distintivas en su composición derivado de fuentes
naturales. Procesos físicos y químicos generan un fraccionamiento
isotópico característico que permite rastrear la fuente que lo generó.
Un ejemplo de ello es la evaporación, donde la fase líquida sufre un
enriquecimiento de isótopos pesados (ej. \textsuperscript{2}H y
\textsuperscript{18}O). Asimismo, han sido ampliamente utilizados para
identificar fuentes de contaminación (Hoefs \& Harmon, 2022a;
Nakajima et al., 2026).

Por otro lado, el incremento de las zonas urbanas es reconocido como
foco de riesgos ambientales en los que los contaminantes históricos
interactúan con la variabilidad climática para generar efectos en
cadena sobre los ecosistemas, la salud pública y la resiliencia urbana
(Li et al., 2026). En México, una de las ciudades más contaminadas es
Salamanca, Guanajuato. La causa principal de la contaminación es la
presencia de industria química y de generación de energía
(Barrón-Adame et al., 2012).

\section{Justificación}
La contaminación atmosférica es un problema ambiental creciente que 
afecta tanto la calidad del aire como el equilibrio del ciclo hidrológico. 
En particular, el vapor de agua atmosférico es clave en la regulación 
del clima y en la distribución de contaminantes, pero su relación con la 
contaminación antropogénica ha sido poco explorada en entornos urbanos y 
semiáridos como Salamanca, Guanajuato.

Mediante el análisis de isótopos estables de oxígeno y deuterio
($\delta^{18}$O y $\delta^{2}$D) del vapor de agua atmosférico, es
posible evaluar su interacción con las emisiones de gases criterio.
Estos estudios proporcionarán una nueva perspectiva sobre los impactos
ambientales de la urbanización y la actividad antropogénica en la
región.

Esta investigación busca comprender la influencia de las actividades 
humanas en la composición isotópica del vapor de agua atmosférico, lo 
que permitirá utilizarlo como un trazador de la contaminación ambiental. 
Los resultados de este estudio pueden contribuir al desarrollo de estrategias 
de monitoreo y mitigación de la contaminación atmosférica.

\section{Hipótesis}
El vapor de agua emitido durante la combustión de combustibles fósiles altera 
la composición isotópica del vapor atmosférico en la ciudad de Salamanca, Guanajuato. 
Por ello, puede utilizarse como un trazador de la intensidad de la contaminación atmosférica
de origen antropogénico.

\section{Objetivo general}


\section{Objetivos específicos}
\begin{itemize}
    \item Identificar las condiciones meteorológicas y las actividades
    antropogénicas que influyen en las variaciones de concentración de
    los contaminantes criterio atmosféricos en la ciudad de Salamanca.
    
    \item Determinar el fraccionamiento isotópico de $\delta^{18}$O y
    $\delta^{2}$D, y su variación temporal en el vapor de agua
    atmosférico de la ciudad de Salamanca.
    
    \item Evaluar el impacto de contaminantes criterio sobre la variación
    del fraccionamiento isotópico de $\delta^{18}$O y $\delta^{2}$D en
    el vapor atmosférico.
    
    \item Generar un modelo de dispersión de contaminantes emitidos por
    las principales industrias de la zona urbana de Salamanca.
\end{itemize}